\documentclass[sigconf,screen,nonacm]{acmart}

% Entfernt den ACM Reference Block und die Zeilennummern (falls vorhanden)
\settopmatter{printacmref=false, printccs=false, printfolios=false}

% Entfernt das Copyright-Datum / Fußzeile unten links ganz unten
\renewcommand\footnotetextcopyrightpermission[1]{}

\begin{document}

\title{PluggedIn: Representing Electric Vehicle Charging Accessibility}

\author{Jungwon Youn}
\affiliation{\institution{University of Victoria}\city{Victoria}\state{BC}\country{Canada}}
\email{jungwonyoun@uvic.ca}

\author{Sunghun Cho}
\affiliation{\institution{University of Victoria}\city{Victoria}\state{BC}\country{Canada}}
\email{shcho980901@uvic.ca}

\author{Elias Dockal}
\affiliation{\institution{University of Victoria}\city{Victoria}\state{BC}\country{Canada}}
\email{eliasdockal@uvic.ca}

\renewcommand{\shortauthors}{Youn, Cho, and Dockal}

\begin{abstract}
EV charging usability depends on multiple factors, including connector compatibility, operating hours, access conditions, and network information. This project compares existing charging infrastructure data, interoperability protocols, and semantic models to identify differences and gaps in how these factors are represented. Based on this comparison, we design PluggedIn, a data model centered on vehicle-specific practical usability. PluggedIn integrates these factors while explicitly representing adapter compatibility, driver network membership, and structured operating periods, providing a foundation for future knowledge graph representation. The project repository, data models, and related artifacts are available at \url{https://github.com/chohoon901/PluggedIn}.
\end{abstract}

\keywords{EV charging, relational data modeling, SQLite, adapter compatibility, data accessibility}

\maketitle

\section{Problem Definition}
Our goal is to let drivers and fleet managers retrieve station availability, connector compatibility, operating hours, and network payment conditions. Flat station and infrastructure lists lose vehicle-to-port suitability associations, adapter requirements, and membership constraints. 

We therefore link each charging station to structured operating periods, individual port blocks, standardized connector types, and organizational roles. Furthermore, we explicitly associate drivers with network memberships and evaluate vehicle charging feasibility through direct connections or adapter-mediated compatibility. 

Our database must answer critical operational questions: \textit{``Can a specific vehicle charge at this station, either directly or via an adapter, and is it currently open and accessible under the driver's network membership?''} State your domain, intended users, and the questions your database must answer. Explain how your model addresses those questions.

\section{Proposed Relational Model}
\subsection{Station, Organization, and Core Entity Relations}
\textbf{Station} records spatial coordinates, address details (street, city, province), truck parking availability, maximum vehicle class restrictions, access restrictions, guest payment permissions, and operational status. \textbf{Organization} describes reusable organizational entities, distinguishing site owners, network operators, and charging networks through an ISA-equivalent relational pattern. Business identifiers and station names serve as mandatory candidate keys. Every relation has an autogenerated integer primary key named \texttt{Id}. The complete relation list appears on the final page.

\subsection{Ports, Operating Periods, and Adapter Relations}
\textbf{OperatingPeriod} records temporal availability by isolating opening and closing times per day, replacing unstructured free-text hours. \textbf{Port} records target-specific outlets and references one station, capturing current types, power capacities, and quantities. \textbf{ConnectorType} standardizes plug nomenclature across disparate sources. To solve real-world connector mismatches, \textbf{Adapter} records hardware specifications, while \textbf{AdapterCompatibility} links adapters to plug-side and inter-side connector types.

\subsection{Vehicles, Drivers, and Membership Relations}
\textbf{Vehicle} stores specifications including VIN, licence plate, make, model, class, range, and maximum charging power. \textbf{Driver} captures user details and commercial carrier affiliations. The many-to-many relationship between vehicles and supported connectors becomes \textbf{VehicleConnector}. Similarly, the many-to-many relationship linking drivers to charging networks becomes \textbf{MemberOf}, capturing validity start/end dates and associated payment methods with a unique pair of foreign keys.

\subsection{Relationship Mapping and Integrity Constraints}
One-to-many relationships use foreign keys on the many side. Business candidate keys use both \texttt{UNIQUE} and \texttt{NOT NULL}. Allowed operational statuses and current types use \texttt{CHECK} constraints where applicable. Foreign-key enforcement must be enabled on every SQLite connection via \texttt{PRAGMA foreign\_keys = ON;}. Referenced records cannot be deleted while dependent rows remain; no cascade is assumed. 
Mapping assumptions. Organizations, stations, vehicles, and drivers may exist before their related transactional rows or memberships are entered; hence parent-side participation is optional. Every port and operating period must reference an existing station. Input validation enforces fixed ISO-style formats for temporal and spatial data.

\section{Relational Operations}
\subsection{Creating and Populating Relations}

\begin{verbatim}

\end{verbatim}

\subsection{SELECT Queries: Adapter-Mediated Compatibility}

\begin{verbatim}

\end{verbatim}

\subsection{INSERT, UPDATE, and DELETE Operations}

\begin{verbatim}

\end{verbatim}

\end{document}